\section{The protocol}
\label{sec:protocol}

\paragraph{Notation and estimand.} Write $C$ for the \textbf{target construct} the benchmark names, $E$ for the \textbf{elicitation intervention} that is actually manipulated, and $V$ for the \textbf{content} the intervention is applied to. The crux of this paper is that in behavioral benchmarks $E$ is typically \emph{not} a valid independent operationalization of $C$: manipulating $E$ establishes that a model was asked in a particular way, not that $C$ obtained.

We therefore ask how detector output responds to $E$ with $V$ held fixed, and to $V$ with $E$ held fixed, as \emph{signed} effects rather than accuracies. In the $2\!\times\!2$ design of §\ref{sec:d2} we fit $Y=\beta_0+\beta_V V+\beta_E E+\beta_{VE}(V\!\times\!E)$ by ordinary least squares, reading $\beta_E$ as the response to the elicitation at fixed content, $\beta_V$ as the response to content at fixed elicitation, and $\beta_{VE}$ as their mutual dependence, with clustered intervals and permutation tests for each. Accuracy is the wrong summary: a detector below 50\% is \emph{anti}-correlated with the label, a statement about direction, not absence of information.

Crucially, \textbf{none of these coefficients alone estimates a causal deception effect}, because deception is not randomized independently of $E$ in this benchmark: $E$ is what produces it. Because $V$ is claim truth-value rather than deception status, $\beta_V$ is likewise not a deception effect: it tests only whether detector output varies systematically with claim content once elicitation is held fixed. What the design does identify is whether the score can be attributed to the construct at all, which is the question we ask.

\paragraph{Diagnostic 1: condition equalization.} Replace the asymmetric elicitation prompts with a single neutral prompt applied to both arms, so that only the construct-bearing content varies. Report the drop. A large drop means the reported score depended on the elicitation asymmetry. This diagnostic is cheap and is the first thing to run, but it is \emph{ambiguous on its own}: removing the instruction may also remove the construct, so a collapse is consistent both with ``the benchmark measured the instruction'' and with ``we deleted the thing being measured.'' Diagnostic 2 exists to resolve that ambiguity.

\paragraph{Diagnostic 2: factorial decomposition.} Cross the construct with the elicitation condition and populate all four cells, including the off-diagonal cell that is usually absent from benchmarks: construct present \emph{without} the usual content cue, and content cue present \emph{without} the construct. Estimate the within-condition contrast. This is the diagnostic that identifies the estimand above; the other two are supporting evidence.

\paragraph{Diagnostic 3: group-held-out evaluation.} Benchmarks reuse a fixed item set across models and conditions. Trial-level cross-validation then leaks: the same claim, or its matched counterpart, can appear in both train and test. Hold out entire semantic groups (here, matched claim pairs) rather than trials, and report the group-clustered interval rather than the trial-level one. If the two agree, memorization is ruled out; if they diverge, the trial-level figure was inflated.

\paragraph{The three are not of a kind.} Diagnostics 1 and 2 address \emph{causal identification}: whether the score can be attributed to the construct. Diagnostic 3 addresses \emph{generalization} (whether the score survives held-out semantic units), which is a different question and can pass while the first two fail. We therefore describe the protocol as two causal-validity diagnostics plus one generalization diagnostic, rather than three of a kind.

\paragraph{Reporting.} We recommend all three be reported alongside any behavioral benchmark score, together with a \emph{surface baseline} (a regex or equivalent applied at a fixed threshold), since a detector that does not clearly exceed one has not demonstrated that it measures anything beyond lexical form. Such a baseline is parameter-free at evaluation time but its patterns are still a researcher choice, so its sensitivity to them should be reported too (§\ref{sec:baseline}). Uncertainty should be clustered at the level of the repeated semantic unit (here, the claim pair): in our data this widens intervals by at most 2\,pp, but that is a fact to be checked rather than assumed.

\paragraph{What the construct is, in this case study.} Our manipulation establishes that a model was \emph{instructed} to deceive; it does not establish an internal intention to deceive. Throughout, therefore, the operationalized construct is \textbf{instruction-conditioned deceptive behavior}. This is not a hedge but the point: we hold ourselves to the standard we apply to prior work, and it sharpens the thesis: the benchmark cannot identify latent deception from this protocol because its purported ground truth is itself entangled with the elicitation manipulation.

\paragraph{Multiplicity.} Confirmatory conclusions rest on the pre-specified aggregate effects (the equalization collapse and the factorial coefficients); we do not apply multiplicity correction to the exploratory per-target analyses, which we label as such.

\paragraph{Case-study setup.} Seven targets (Llama~3.2~3B, Llama~3.1~8B, Llama~3.3~70B~\citep{meta2024llama3}, Mistral~7B~\citep{jiang2023mistral}, Qwen~2.5~7B/14B~\citep{qwen2025qwen25}, Claude Haiku~4.5~\citep{anthropic2024claude}), $n=93$--100 balanced trials each (689 total); Qwen~2.5~32B additionally in the replication of \citet{pacchiardi2023catch}. Claims are 50 matched true/false factual pairs (``Water boils at 100\textdegree C'' vs.\ ``at 85\textdegree C''), giving \textbf{50 independent semantic units}, the level at which we cluster every uncertainty estimate. The \emph{instructed} condition uses asymmetric truth/lie system prompts; the \emph{equalized} condition gives both arms an identical neutral prompt so that only claim truth-value varies. Full protocol in Appendix~\ref{app:setup}.
